% ============================================================================
% ONR FY2027 YIP (NOFO N0001426SF004): Budget Justification
% Attach under Section L of the R&R Budget form (NOFO p. 21). No page limit.
% Numbers come from budget.py (budget-table.tex, budget-items.tex); edit the
% assumptions there and rebuild, rather than editing numbers here.
% Build: python budget.py && xelatex budget-justification.tex (twice)
% ============================================================================
\documentclass[11pt,letterpaper]{article}
\usepackage[margin=1in]{geometry}
\usepackage{fontspec}
\IfFontExistsTF{Times New Roman}{\setmainfont{Times New Roman}}{\setmainfont{TeX Gyre Termes}}
\usepackage{microtype}
\usepackage{xcolor}
\usepackage{array}
\usepackage{booktabs}
\usepackage{longtable}
\usepackage[colorlinks=true,urlcolor=blue!70!black,linkcolor=blue!70!black]{hyperref}
\setlength{\parindent}{0pt}
\setlength{\parskip}{5pt}
\newcommand{\todo}[1]{\textcolor{red}{[TODO Sterling or BYU RAO: #1]}}
\pagestyle{plain}
\newcommand{\BudgetTotal}{747,651}
\newcommand{\BudgetYearOne}{248,959}
\newcommand{\MTDCYearOne}{140,298}
\newcommand{\FAYearOne}{72,253}
\newcommand{\BuildsYearOne}{28}
\newcommand{\BudgetYearTwo}{248,828}
\newcommand{\MTDCYearTwo}{151,728}
\newcommand{\FAYearTwo}{78,140}
\newcommand{\BuildsYearTwo}{29}
\newcommand{\BudgetYearThree}{249,864}
\newcommand{\MTDCYearThree}{152,036}
\newcommand{\FAYearThree}{78,299}
\newcommand{\BuildsYearThree}{26}


\begin{document}

\begin{center}
{\large\bfseries Budget Justification\par}
\vspace{4pt}
Office of Naval Research FY2027 Young Investigator Program, NOFO N0001426SF004\\
Cost-Aware Selection of Experiments for Printable, Corrosion-Resistant
High-Strength Aluminum Alloys\\
Principal Investigator: Sterling G.\ Baird, Brigham Young University (BYU)\\
Period of performance: 1 April 2027 to 31 March 2030 (36 months, base period only)
\end{center}

\textcolor{red}{Draft of 10 October 2026. Red items are assumptions that
BYU Research Administration must replace with current figures before the
budget is entered in Kuali.}

\section*{Summary}

The request is \$\BudgetTotal{} over 36 months, below the \$250,000 limit in each
12-month period (NOFO p.~21). All amounts are rounded to the nearest dollar.
There are no subawards, consultants, participant support costs, or option
periods.

\begin{center}
\begin{tabular}{@{}lrrrr@{}}
\toprule
\textbf{Category} & \textbf{Year 1} & \textbf{Year 2} & \textbf{Year 3} & \textbf{Total} \\
\midrule
A--C. Personnel and fringe benefits (pooled) & 87,648 & 90,278 & 92,986 & 270,912 \\
%
%
D. Equipment & 18,000 & 0 & 0 & 18,000 \\
E. Travel & 6,600 & 8,600 & 8,600 & 23,800 \\
F. Materials and supplies & 15,350 & 15,350 & 15,350 & 46,050 \\
G. Other direct costs & 49,108 & 56,460 & 54,629 & 160,197 \\
\midrule
Total direct costs & 176,706 & 170,688 & 171,565 & 518,959 \\
Modified total direct costs (MTDC) & 140,298 & 151,728 & 152,036 & 444,062 \\
H. F\&A (51.5\% of MTDC) & 72,253 & 78,140 & 78,299 & 228,692 \\
\midrule
\textbf{Total} & \textbf{248,959} & \textbf{248,828} & \textbf{249,864} & \textbf{747,651} \\
\bottomrule
\end{tabular}

\end{center}

\section*{Rates and bases}

\begin{itemize}\setlength{\itemsep}{1pt}
\item \textbf{Escalation.} Salaries, wages, and tuition rise 3\% per year
after Year 1, the inflation rate BYU's budget guide suggests.
\item \textbf{Fringe benefits.} From BYU's 2026 budget planning sheet: 19.7\%
of faculty spring and summer salary; 0\% for students while enrolled; 7.7\%
(FICA) for students when not enrolled, taken here as the four summer months
of student wages \todo{confirm the 2027 rates, which Kuali fills in, and
whether the graduate students will be enrolled in spring and summer}.
\item \textbf{Facilities and administrative (F\&A) costs.} 51.5\% of modified
total direct costs (MTDC), BYU's on-campus research rate negotiated with the
Department of Health and Human Services (agreement of 2 August 2023, with
the same rate provisional from 1 January 2025 until amended)
\todo{attach the current agreement from BYU's Sponsored Projects Office, as
the NOFO requires (p.~21)}. MTDC excludes equipment and tuition remission.
\item \textbf{Single audit.} \todo{insert the link to BYU's most recent
single audit report in the Federal Audit Clearinghouse (NOFO p.~21 asks for
it); BYU's UEI is JWSYC7RUMJD1.}
\end{itemize}

\section*{A. Senior personnel}

Sterling G.\ Baird, Principal Investigator, summer salary only [per-person salary detail removed].
%
%
%
%
%
The PI leads the method development, sets the campaign
plans, and supervises the students.

\section*{B. Other personnel}

\textbf{Two Ph.D.\ research assistants}, 12 months each for 3 years
[per-person rate removed].
Student 1 develops the multi-objective and
asynchronous search rules, the benchmarks, and the software (Task 1).
Student 2 runs the experimental platform: powder, builds, microscopy,
tensile tests, and corrosion tests (Task 2). Both run the live campaigns
(Task 3).

\textbf{One undergraduate researcher}, 500 hours per year [per-person rate removed]
\todo{confirm BYU's student wage scale}, who weighs feedstock, operates the
atomizer under supervision, prepares metallographic specimens, and runs the
electrochemical corrosion screen.

\section*{C. Fringe benefits}

Faculty fringe is 19.7\% of the PI's summer salary. Student fringe is 7.7\%
of the wages paid in the four summer months, when students are assumed not
to be enrolled, and 0\% otherwise. Year-by-year amounts are in the tables
below.

\section*{D. Equipment}

\textbf{Potentiostat with corrosion software and a flat cell} (Year 1,
\$18,000 estimate) \todo{obtain a vendor quote, for example for a Gamry
Interface 1010E or BioLogic SP-150e with a flat cell; the NOFO requires
quotes for equipment (p.~22)}. The electrochemical corrosion screen is one
step of the workflow that the search schedules (Technical Proposal,
Section~1), so the instrument must be available every day and controlled by
the lab's own scripts. A shared instrument booked by the hour cannot be
scheduled by the search this way \todo{confirm that no BYU potentiostat can
be dedicated to the project; if one can, delete this line, and Year 1 has room
for about 14 more builds}. Equipment is excluded from MTDC.

\section*{E. Travel}

Each trip is estimated from airfare, lodging and meals at General Services
Administration (GSA) rates, ground transport, and registration where it
applies \todo{rebuild each trip on ONR's travel estimate spreadsheet with
current GSA rates}.
\begin{itemize}\setlength{\itemsep}{1pt}
\item \textbf{ONR program review} (each year): Provo, UT to Arlington, VA;
PI; 3 days; \$2,000. The NOFO allows PI travel to ONR peer reviews
(pp.~22, 38).
\item \textbf{Conference} (each year): TMS Annual Meeting or Materials Science
\& Technology; PI and one student; 4 days; \$2,300 per traveler including
registration; \$4,600. Results are presented to the materials and
manufacturing community.
\item \textbf{Navy laboratory visit} (Years 2 and 3): Provo, UT to NSWC
Carderock Division, West Bethesda, MD; PI; 3 days; \$2,000. Purpose: share
results and corrosion data with Navy materials engineers and discuss
transition.
\end{itemize}

\section*{F. Materials and supplies}

Itemized with quantities and unit prices in the tables below. The basis is
current catalog prices \todo{obtain quotes for master alloys, especially the
aluminum-scandium master alloy, and confirm the argon use per atomizer run}.
Feedstock is weighed into machined crucibles of high-purity aluminum, about
100~g per composition, and atomized in the PI's laboratory.

\section*{G. Other direct costs}

\begin{itemize}\setlength{\itemsep}{1pt}
\item \textbf{LPBF builds}, University of Utah user facility, \$800 per build,
bought on a purchase order (not a subaward). The facility's staff run each
build to the laser settings that the search chooses. The number of builds
(\BuildsYearOne, \BuildsYearTwo, and \BuildsYearThree) is what remains
under the annual limit after the other costs; builds are the step that sets
the pace of the campaigns.
\item \textbf{BYU Microscopy Center} (SEM, EDS, EBSD), 100 hours per year at
\$45 per hour \todo{confirm the internal rates}.
\item \textbf{ICP-MS composition checks} at BYU, 40 samples per year at \$35
\todo{confirm the internal rate}.
\item \textbf{Tensile specimen machining} at BYU, 60 specimens per year at \$40
\todo{confirm the shop rate}.
\item \textbf{Neutral salt-spray exposure} (ASTM B117) at an outside test lab
for the leading candidates, two runs per year in Years 2 and 3 at \$1,500
\todo{obtain a quote, or use a BYU chamber if one exists}.
\item \textbf{Open-access publication charges}, one article per year in Years
2 and 3, \$3,000 each.
\item \textbf{Graduate tuition} for the two Ph.D.\ students, fall and winter
semesters. BYU's 2026--27 graduate tuition is \$4,468 per semester for
Latter-day Saint students and \$8,936 for others; the budget assumes the
first rate plus 3\% for 2027--28 \todo{double this line for a student who
pays the non-member rate}. BYU charges tuition to a sponsored project only
with a research assistant wage, as its own line. Excluded from MTDC.
\end{itemize}

\section*{H. Indirect costs}

F\&A is 51.5\% of MTDC. MTDC is total direct costs less equipment and tuition:
\$\MTDCYearOne{} in Year 1, \$\MTDCYearTwo{} in Year 2, and \$\MTDCYearThree{}
in Year 3, for F\&A of \$\FAYearOne, \$\FAYearTwo, and \$\FAYearThree.

\section*{Itemized budget by year}

\subsection*{Year 1 (1 Apr 2027 to 31 Mar 2028)}
\begin{longtable}{@{}>{\raggedright\arraybackslash}p{2.6in}>{\raggedright\arraybackslash}p{2.75in}r@{}}
\toprule \textbf{Item} & \textbf{Basis} & \textbf{\$} \\ \midrule \endhead
\multicolumn{3}{@{}l}{\textit{A--C. Personnel and fringe benefits (pooled)}} \\
PI, two Ph.D. research assistants, one undergraduate researcher, and fringe benefits & pooled; per-person detail removed & 87,648 \\
%
%
%
%
%
%
\multicolumn{3}{@{}l}{\textit{D. Equipment}} \\
Potentiostat with corrosion software and flat cell & vendor quote needed (\textcolor{red}{TODO}) & 18,000 \\
\multicolumn{3}{@{}l}{\textit{E. Travel}} \\
ONR program review, Arlington VA (PI, 3 days) & per-trip estimate; \textcolor{red}{TODO} GSA rates & 2,000 \\
TMS Annual Meeting or MS\&T (PI and one student, 4 days) & per-trip estimate; \textcolor{red}{TODO} GSA rates & 4,600 \\
\multicolumn{3}{@{}l}{\textit{F. Materials and supplies}} \\
Master alloys and elemental powders (Al-Zr, Al-Sc, Al-Mn, Al-Cr, Mg, Zn, Cu) & catalog estimate; \textcolor{red}{TODO} quotes & 6,000 \\
High-purity aluminum bar for machined crucibles (30 kg at \$45/kg) & catalog estimate; \textcolor{red}{TODO} quotes & 1,350 \\
Argon for atomization and powder handling (30 cylinders at \$90) & catalog estimate; \textcolor{red}{TODO} quotes & 2,700 \\
Metallography consumables (papers, cloths, suspensions, mounting) & catalog estimate; \textcolor{red}{TODO} quotes & 3,000 \\
Corrosion consumables (NaCl, H2O2, HNO3, reference electrodes, gaskets) & catalog estimate; \textcolor{red}{TODO} quotes & 1,500 \\
Sieves, powder containers, and inert-gas glove bags & catalog estimate; \textcolor{red}{TODO} quotes & 800 \\
\multicolumn{3}{@{}l}{\textit{G. Other direct costs}} \\
LPBF builds, University of Utah user facility (28) & 28 $\times$ \$800 per build (purchase order) & 22,400 \\
BYU Microscopy Center (SEM, EDS, EBSD) & 100 h $\times$ \$45/h; \textcolor{red}{TODO} internal rates & 4,500 \\
ICP-MS composition checks (BYU) & 40 samples $\times$ \$35; \textcolor{red}{TODO} internal rate & 1,400 \\
Tensile specimen machining (BYU) & 60 specimens $\times$ \$40; \textcolor{red}{TODO} shop rate & 2,400 \\
Graduate tuition (2 students) & 2 $\times$ \$9,204; excluded from MTDC & 18,408 \\
\midrule
Total direct costs & & 176,706 \\
MTDC (direct costs less equipment and tuition) & & 140,298 \\
H. F\&A at 51.5\% of MTDC & & 72,253 \\
\textbf{Year 1 total} & & \textbf{248,959} \\
\bottomrule
\end{longtable}
\subsection*{Year 2 (1 Apr 2028 to 31 Mar 2029)}
\begin{longtable}{@{}>{\raggedright\arraybackslash}p{2.6in}>{\raggedright\arraybackslash}p{2.75in}r@{}}
\toprule \textbf{Item} & \textbf{Basis} & \textbf{\$} \\ \midrule \endhead
\multicolumn{3}{@{}l}{\textit{A--C. Personnel and fringe benefits (pooled)}} \\
PI, two Ph.D. research assistants, one undergraduate researcher, and fringe benefits & pooled; per-person detail removed & 90,278 \\
%
%
%
%
%
%
\multicolumn{3}{@{}l}{\textit{E. Travel}} \\
ONR program review, Arlington VA (PI, 3 days) & per-trip estimate; \textcolor{red}{TODO} GSA rates & 2,000 \\
TMS Annual Meeting or MS\&T (PI and one student, 4 days) & per-trip estimate; \textcolor{red}{TODO} GSA rates & 4,600 \\
NSWC Carderock, West Bethesda MD (PI, 3 days) & per-trip estimate; \textcolor{red}{TODO} GSA rates & 2,000 \\
\multicolumn{3}{@{}l}{\textit{F. Materials and supplies}} \\
Master alloys and elemental powders (Al-Zr, Al-Sc, Al-Mn, Al-Cr, Mg, Zn, Cu) & catalog estimate; \textcolor{red}{TODO} quotes & 6,000 \\
High-purity aluminum bar for machined crucibles (30 kg at \$45/kg) & catalog estimate; \textcolor{red}{TODO} quotes & 1,350 \\
Argon for atomization and powder handling (30 cylinders at \$90) & catalog estimate; \textcolor{red}{TODO} quotes & 2,700 \\
Metallography consumables (papers, cloths, suspensions, mounting) & catalog estimate; \textcolor{red}{TODO} quotes & 3,000 \\
Corrosion consumables (NaCl, H2O2, HNO3, reference electrodes, gaskets) & catalog estimate; \textcolor{red}{TODO} quotes & 1,500 \\
Sieves, powder containers, and inert-gas glove bags & catalog estimate; \textcolor{red}{TODO} quotes & 800 \\
\multicolumn{3}{@{}l}{\textit{G. Other direct costs}} \\
LPBF builds, University of Utah user facility (29) & 29 $\times$ \$800 per build (purchase order) & 23,200 \\
BYU Microscopy Center (SEM, EDS, EBSD) & 100 h $\times$ \$45/h; \textcolor{red}{TODO} internal rates & 4,500 \\
ICP-MS composition checks (BYU) & 40 samples $\times$ \$35; \textcolor{red}{TODO} internal rate & 1,400 \\
Tensile specimen machining (BYU) & 60 specimens $\times$ \$40; \textcolor{red}{TODO} shop rate & 2,400 \\
Neutral salt-spray exposure (ASTM B117), outside lab & 2 runs $\times$ \$1,500; \textcolor{red}{TODO} quote & 3,000 \\
Open-access publication charges & one article per year & 3,000 \\
Graduate tuition (2 students) & 2 $\times$ \$9,480; excluded from MTDC & 18,960 \\
\midrule
Total direct costs & & 170,688 \\
MTDC (direct costs less equipment and tuition) & & 151,728 \\
H. F\&A at 51.5\% of MTDC & & 78,140 \\
\textbf{Year 2 total} & & \textbf{248,828} \\
\bottomrule
\end{longtable}
\subsection*{Year 3 (1 Apr 2029 to 31 Mar 2030)}
\begin{longtable}{@{}>{\raggedright\arraybackslash}p{2.6in}>{\raggedright\arraybackslash}p{2.75in}r@{}}
\toprule \textbf{Item} & \textbf{Basis} & \textbf{\$} \\ \midrule \endhead
\multicolumn{3}{@{}l}{\textit{A--C. Personnel and fringe benefits (pooled)}} \\
PI, two Ph.D. research assistants, one undergraduate researcher, and fringe benefits & pooled; per-person detail removed & 92,986 \\
%
%
%
%
%
%
\multicolumn{3}{@{}l}{\textit{E. Travel}} \\
ONR program review, Arlington VA (PI, 3 days) & per-trip estimate; \textcolor{red}{TODO} GSA rates & 2,000 \\
TMS Annual Meeting or MS\&T (PI and one student, 4 days) & per-trip estimate; \textcolor{red}{TODO} GSA rates & 4,600 \\
NSWC Carderock, West Bethesda MD (PI, 3 days) & per-trip estimate; \textcolor{red}{TODO} GSA rates & 2,000 \\
\multicolumn{3}{@{}l}{\textit{F. Materials and supplies}} \\
Master alloys and elemental powders (Al-Zr, Al-Sc, Al-Mn, Al-Cr, Mg, Zn, Cu) & catalog estimate; \textcolor{red}{TODO} quotes & 6,000 \\
High-purity aluminum bar for machined crucibles (30 kg at \$45/kg) & catalog estimate; \textcolor{red}{TODO} quotes & 1,350 \\
Argon for atomization and powder handling (30 cylinders at \$90) & catalog estimate; \textcolor{red}{TODO} quotes & 2,700 \\
Metallography consumables (papers, cloths, suspensions, mounting) & catalog estimate; \textcolor{red}{TODO} quotes & 3,000 \\
Corrosion consumables (NaCl, H2O2, HNO3, reference electrodes, gaskets) & catalog estimate; \textcolor{red}{TODO} quotes & 1,500 \\
Sieves, powder containers, and inert-gas glove bags & catalog estimate; \textcolor{red}{TODO} quotes & 800 \\
\multicolumn{3}{@{}l}{\textit{G. Other direct costs}} \\
LPBF builds, University of Utah user facility (26) & 26 $\times$ \$800 per build (purchase order) & 20,800 \\
BYU Microscopy Center (SEM, EDS, EBSD) & 100 h $\times$ \$45/h; \textcolor{red}{TODO} internal rates & 4,500 \\
ICP-MS composition checks (BYU) & 40 samples $\times$ \$35; \textcolor{red}{TODO} internal rate & 1,400 \\
Tensile specimen machining (BYU) & 60 specimens $\times$ \$40; \textcolor{red}{TODO} shop rate & 2,400 \\
Neutral salt-spray exposure (ASTM B117), outside lab & 2 runs $\times$ \$1,500; \textcolor{red}{TODO} quote & 3,000 \\
Open-access publication charges & one article per year & 3,000 \\
Graduate tuition (2 students) & 2 $\times$ \$9,765; excluded from MTDC & 19,529 \\
\midrule
Total direct costs & & 171,565 \\
MTDC (direct costs less equipment and tuition) & & 152,036 \\
H. F\&A at 51.5\% of MTDC & & 78,299 \\
\textbf{Year 3 total} & & \textbf{249,864} \\
\bottomrule
\end{longtable}


\end{document}
